\par\Needspace{6\baselineskip}\medskip
\hypertarget{native-076-article058}{}\pdfbookmark[2]{58.}{article-058}
\begin{center}58.\end{center}
Our formulas for the distances from given planes involve $v$ and $r$. Whenever these quantities have to be determined first from the time, some of the operations can be shortened further, appreciably reducing the work. A very simple formula gives the distances directly from the eccentric anomaly in the ellipse, or from the auxiliary quantity $F$ or $u$ in the hyperbola. This entirely dispenses with calculating the true anomaly and the radius vector. The expression $kr\sin(v+K)$ is transformed as follows.

I. \textit{For the ellipse}, retaining the notation of article 8, it becomes
\[
ak\cos\varphi\cos K\sin E+ak\sin K(\cos E-e).
\]
Thus, if we determine $l$, $L$, $\lambda$ by the equations
\[
\begin{aligned}
ak\sin K&=l\sin L\\
ak\cos\varphi\cos K&=l\cos L\\
-eak\sin K&=-el\sin L=\lambda,
\end{aligned}
\]
\SourceLeaf{077}our expression becomes $l\sin(E+L)+\lambda$, where $l$, $L$, $\lambda$ will be constant insofar as $k$, $K$, $e$ can be regarded as constant. Otherwise the remarks in the preceding article apply equally to the calculation of their changes.

As an example, we give the transformation of the expression for $x$ found in article 56, setting the longitude of perihelion $=121^0\,17^{\prime}\,34^{\prime\prime}{,}4$, $\varphi=14^0\,13^{\prime}\,31^{\prime\prime}{,}97$, and $\log a=0{,}442\,3790$. The distance of perihelion from the ascending node in the ecliptic is thus $=308^0\,49^{\prime}\,20^{\prime\prime}{,}7=u-v$. Hence $K=212^0\,36^{\prime}\,56^{\prime\prime}{,}09$. We therefore have
\begin{center}
\begin{tabular}{@{}ll@{}}
\makebox[13em][l]{$\log ak$\hfill .\,.\,.\,.\,.\,.\,.} & $0{,}441\,1713$\\[.3em]
\makebox[13em][l]{$\log\sin K$\hfill .\,.\,.\,.\,.} & $9{,}731\,5887_n$\\[.3em]
\makebox[13em][l]{$\log ak\cos\varphi$\hfill .\,.\,.} & $0{,}427\,6456$\\[.3em]
\makebox[13em][l]{$\log\cos K$\hfill .\,.\,.\,.\,.} & $9{,}925\,4698_n$
\end{tabular}
\end{center}
\begin{center}
\begin{tabular}{@{}ll@{}}
\makebox[13em][l]{$\log l\sin L$\hfill .\,.\,.\,.} & $0{,}172\,7600_n$\\[.3em]
\makebox[13em][l]{$\log l\cos L$\hfill .\,.\,.\,.} & $0{,}353\,1154_n$\\[.3em]
\hline
hence $L=$ & $213^0\,25^{\prime}\,51^{\prime\prime}{,}31$\\[.3em]
\makebox[13em][l]{$\log l$\hfill .\,.\,.\,.\,.\,.\,.\,.} & $0{,}431\,6627$\\[.3em]
\makebox[13em][l]{$\log\lambda$\hfill .\,.\,.\,.\,.\,.\,.\,.} & $9{,}563\,2352$\\[.3em]
$\lambda=$ & $+0{,}365\,7929$
\end{tabular}
\end{center}

II. In the hyperbola, the formula $kr\sin(v+K)$ becomes, by article 21,
\[
\lambda+\mu\operatorname{tang}F+\nu\operatorname{secans}F,
\]
if we set
\[
ebk\sin K=\lambda,\qquad bk\operatorname{tang}\psi\cos K=\mu,\qquad -bk\sin K=\nu;
\]
the same expression can evidently also be reduced to the form $\frac{n\sin(F+N)+\nu}{\cos F}$. If, in place of $F$, we use the auxiliary quantity $u$, the expression $kr\sin(v+K)$ becomes $\alpha+\beta u+\frac{\gamma}{u}$ by article 21, where $\alpha$, $\beta$, $\gamma$ are determined by the formulas
\[
\begin{aligned}
\alpha&=\lambda=ebk\sin K\\
\beta&=\tfrac{1}{2}(\nu+\mu)=-\tfrac{1}{2}ebk\sin(K-\psi)\\
\gamma&=\tfrac{1}{2}(\nu-\mu)=-\tfrac{1}{2}ebk\sin(K+\psi).
\end{aligned}
\]

III. In the parabola, where the true anomaly is derived directly from the time, all that remains is to substitute the value of the radius vector. If we denote the distance at perihelion by $q$, the expression $kr\sin(v+K)$ becomes $=\frac{qk\sin(v+K)}{\cos\tfrac{1}{2}v^2}\mathord{\cdot}$
